\lysh{34309}{4}

\begin{alist}
\item Για $\alpha = 1$ ο τύπος της συνάρτησης $g$ γίνεται: $g(x) = x + 1, \ \ x \in \mathbb{R}$.
Οι συναρτήσεις $f$ και $g$ έχουν πεδίο ορισμού το $\mathbb{R}$. Οι τετμημένες των σημείων τομής των
γραφικών τους παραστάσεων είναι οι λύσεις της εξίσωσης:
\begin{align*}
f(x) = g(x) \Leftrightarrow x^2 + 1 = x + 1 \ &\Leftrightarrow \\
\Leftrightarrow x^2 - x = 0 \Leftrightarrow x(x - 1) = 0 \ &\Leftrightarrow \\
x = 0 \text{ ή } x - 1 = 0 \ &\Leftrightarrow \\
x &= 0 \text{ ή } x = 1
\end{align*}
Επίσης, $g(0) = 0 + 1 = 1$ και $g(1) = 1 + 1 = 2$.
Άρα τα σημεία τομής είναι τα $A(0, g(0))$ και $B(1, g(1))$ δηλαδή τα $A(0, 1)$ και $B(1, 2)$.

\item Οι γραφικές παραστάσεις των συναρτήσεων $f$ και $g$ τέμνονται σε δύο σημεία αν και
μόνο αν η εξίσωση
\[f(x) = g(x) \Leftrightarrow x^2 + 1 = x + \alpha\]
έχει δύο άνισες λύσεις (που θα είναι οι τετμημένες $x_1, x_2$ των σημείων αυτών).
Δηλαδή αν και μόνο αν η εξίσωση
$x^2 - x + 1 - \alpha = 0$ (1)
έχει δύο άνισες ρίζες.
Η διακρίνουσα του τριωνύμου $x^2 - x + 1 - \alpha$ είναι
\[\Delta = (-1)^2 - 4 \cdot 1 \cdot (1 - \alpha) = 1 - 4 + 4\alpha = 4\alpha - 3.\]
Η εξίσωση (1) έχει δύο άνισες ρίζες αν και μόνο αν
\begin{align*}
\Delta > 0 \Leftrightarrow 4\alpha - 3 > 0 \ &\Leftrightarrow \\
4\alpha &> 3 \Leftrightarrow \alpha > \dfrac{3}{4}
\end{align*}

\item Επειδή $\alpha > 1 > \dfrac{3}{4}$, η εξίσωση (1) έχει δύο ρίζες άνισες, τις $x_1, x_2$. Το γινόμενό τους είναι:
\[P = x_1 x_2 = \dfrac{\gamma}{\alpha} = 1 - \alpha.\]
Αλλά, $\alpha > 1 \Rightarrow 1 - \alpha < 0$, δηλαδή $P < 0$. Οπότε οι τετμημένες των σημείων τομής των
γραφικών παραστάσεων των συναρτήσεων $f$ και $g$ είναι ετερόσημες.
\end{alist}

